% ===== PRÉAMBULE CANONIQUE QCM — CHIMIE =====
% Sert à compiler controle.tex ET à la revalidation automatique du JSON
% (valider_qcm.py). NE PAS MODIFIER : noms et packages figés.
\documentclass[11pt,a4paper]{article}
\usepackage[utf8]{inputenc}\usepackage[T1]{fontenc}\usepackage{lmodern}
\usepackage[french]{babel}
\usepackage{amsmath,amssymb,mathtools}\usepackage{icomma}
\usepackage[version=4]{mhchem}          % \ce{...} : équations & formules
\usepackage{siunitx}\sisetup{locale=FR} % \qty{}{}, \si{}, \ang{}
\usepackage{chemfig}                    % \chemfig{...} : structures développées
\usepackage{xcolor}\usepackage{colortbl}
\usepackage{tikz}\usepackage{pgfplots}\pgfplotsset{compat=1.18}
\usepackage{enumitem}\usepackage{array,booktabs}
\usepackage[a4paper,margin=2.2cm]{geometry}
\usetikzlibrary{arrows.meta,calc,angles,quotes,positioning,patterns,backgrounds,shapes.geometric}
\newcommand{\R}{\mathbb{R}}\newcommand{\N}{\mathbb{N}}\newcommand{\Z}{\mathbb{Z}}\newcommand{\ds}{\displaystyle}
% ============================================
\begin{document}
\begin{center}\Large\bfseries QCM type concours --- Chimie, Fiche 11 : transfert d'électrons et piles\end{center}
\vspace{4pt}\hrule\vspace{10pt}
\noindent\textbf{CHIM-F11-Q01 --- Nombre d'oxydation d'un atome dans un ion polyatomique}\par
Quel est le nombre d'oxydation du manganèse dans l'ion permanganate \ce{MnO4^-} ? On rappelle que $\mathrm{n.o.}(\ce{O})=-\mathrm{II}$ et que la somme des nombres d'oxydation d'une espèce est égale à sa charge globale.
\begin{enumerate}[label=\Alph*)]
\item $+\mathrm{III}$
\item $+\mathrm{VI}$
\item $+\mathrm{VII}$
\item $+\mathrm{VIII}$
\end{enumerate}
\textit{Bonne réponse : C}\par
L'ion \ce{MnO4^-} porte la charge $-1$. En posant $x=\mathrm{n.o.}(\ce{Mn})$ et $\mathrm{n.o.}(\ce{O})=-\mathrm{II}$ pour les quatre oxygènes : $x+4\times(-2)=-1$, soit $x-8=-1$, d'où $x=+7=+\mathrm{VII}$. C'est le degré d'oxydation maximal du manganèse (élément du groupe 7).\par
\textit{Pièges :}
\begin{itemize}
\item A: règle des peroxydes appliquée à tort, $\mathrm{n.o.}(\ce{O})=-\mathrm{I}$ : $x-4=-1$ donne $+\mathrm{III}$.
\item B: charge de l'ion prise à $-2$ (confusion avec \ce{SO4^{2-}}) : $x-8=-2$ donne $+\mathrm{VI}$.
\item D: charge de l'ion oubliée (somme mise à $0$) : $x-8=0$ donne $+\mathrm{VIII}$, degré pourtant impossible pour \ce{Mn}.
\end{itemize}
\vspace{8pt}\hrule\vspace{8pt}
\noindent\textbf{CHIM-F11-Q02 --- Reconnaître une réaction d'oxydo-réduction}\par
Parmi les réactions suivantes (toutes équilibrées), laquelle n'est \textbf{pas} une réaction d'oxydo-réduction ?
\begin{enumerate}[label=\Alph*)]
\item \ce{CaO + CO2 -> CaCO3}
\item \ce{2 H2 + O2 -> 2 H2O}
\item \ce{Cl2 + H2O -> HCl + HClO}
\item \ce{Cl2 + 2 KBr -> 2 KCl + Br2}
\end{enumerate}
\textit{Bonne réponse : A}\par
Une réaction est redox si et seulement si au moins un élément change de nombre d'oxydation. Dans \ce{CaO + CO2 -> CaCO3}, on a $\mathrm{n.o.}(\ce{Ca})=+\mathrm{II}$, $\mathrm{n.o.}(\ce{C})=+\mathrm{IV}$ et $\mathrm{n.o.}(\ce{O})=-\mathrm{II}$ des deux côtés : aucun degré ne varie. C'est une réaction acide-base entre oxydes (combinaison), sans transfert d'électrons.\par
\textit{Pièges :}
\begin{itemize}
\item B: synthèse à partir des corps simples \ce{H2} et \ce{O2}, prise pour « non redox » alors que \ce{H} passe de $0$ à $+\mathrm{I}$ et \ce{O} de $0$ à $-\mathrm{II}$.
\item C: dismutation du chlore prise pour une simple dissolution ; \ce{Cl} passe simultanément de $0$ à $-\mathrm{I}$ (dans \ce{HCl}) et à $+\mathrm{I}$ (dans \ce{HClO}).
\item D: réaction vue comme un simple échange d'ions, en oubliant que \ce{Cl} passe de $0$ à $-\mathrm{I}$ et \ce{Br} de $-\mathrm{I}$ à $0$.
\end{itemize}
\vspace{8pt}\hrule\vspace{8pt}
\noindent\textbf{CHIM-F11-Q03 --- Équilibrage redox : somme des coefficients stœchiométriques}\par
En milieu acide, l'ion permanganate oxyde l'ion fer(II) : \ce{MnO4^- + Fe^{2+} + H+ -> Mn^{2+} + Fe^{3+} + H2O}. Une fois l'équation correctement équilibrée, que vaut la somme de tous les coefficients stœchiométriques (les « 1 » compris) ?
\begin{enumerate}[label=\Alph*)]
\item $16$
\item $20$
\item $22$
\item $24$
\end{enumerate}
\textit{Bonne réponse : D}\par
Demi-équations : \ce{MnO4^- + 8 H+ + 5 e^- -> Mn^{2+} + 4 H2O} ($\mathrm{n.o.}(\ce{Mn})$ : $+\mathrm{VII}\to+\mathrm{II}$, 5 électrons) et \ce{Fe^{2+} -> Fe^{3+} + e^-}. On multiplie la seconde par 5, d'où \ce{MnO4^- + 5 Fe^{2+} + 8 H+ -> Mn^{2+} + 5 Fe^{3+} + 4 H2O}. Somme $=1+5+8+1+5+4=24$.\par
\textit{Pièges :}
\begin{itemize}
\item A: électrons non égalisés, coefficient $1$ devant \ce{Fe^{2+}}/\ce{Fe^{3+}} au lieu de $5$ : $1+1+8+1+1+4=16$.
\item B: nombre de \ce{H+} pris égal au nombre de \ce{H2O} ($4$) au lieu du double ($8$) : $1+5+4+1+5+4=20$.
\item C: seuls les coefficients explicites (supérieurs à $1$) comptés, les « 1 » implicites oubliés : $5+8+5+4=22$.
\end{itemize}
\vspace{8pt}\hrule\vspace{8pt}
\noindent\textbf{CHIM-F11-Q04 --- Équilibrage : nombre d'électrons d'une demi-équation (dichromate)}\par
En milieu acide, l'ion dichromate \ce{Cr2O7^{2-}} est réduit en ion chrome(III) \ce{Cr^{3+}}. Combien d'électrons sont échangés dans cette demi-équation ?
\begin{enumerate}[label=\Alph*)]
\item $3$
\item $4$
\item $6$
\item $8$
\end{enumerate}
\textit{Bonne réponse : C}\par
Dans \ce{Cr2O7^{2-}}, $2x+7\times(-2)=-2$ donne $x=+\mathrm{VI}$ ; dans \ce{Cr^{3+}}, $\mathrm{n.o.}=+\mathrm{III}$. Chaque atome de chrome gagne $6-3=3$ électrons, mais le dichromate en contient \textbf{deux} : $2\times3=6$ électrons. La demi-équation est \ce{Cr2O7^{2-} + 14 H+ + 6 e^- -> 2 Cr^{3+} + 7 H2O}.\par
\textit{Pièges :}
\begin{itemize}
\item A: $\Delta\mathrm{n.o.}=3$ calculé pour un seul \ce{Cr}, sans multiplier par les deux atomes du dichromate.
\item B: $\mathrm{n.o.}(\ce{Cr})$ pris à $+\mathrm{VII}$ (charge $-2$ de l'ion oubliée) puis $\Delta=7-3=4$, sans le facteur 2.
\item D: même erreur $\mathrm{n.o.}(\ce{Cr})=+\mathrm{VII}$ mais avec le facteur 2 : $(7-3)\times2=8$.
\end{itemize}
\vspace{8pt}\hrule\vspace{8pt}
\noindent\textbf{CHIM-F11-Q05 --- Équilibrage d'une demi-équation en milieu basique}\par
En milieu \textbf{basique}, l'ion permanganate \ce{MnO4^-} est réduit en dioxyde de manganèse \ce{MnO2}. Parmi les demi-équations suivantes, laquelle est correctement équilibrée en masse \emph{et} en charge, dans ce milieu ?
\begin{enumerate}[label=\Alph*)]
\item \ce{MnO4^- + 4 H+ + 3 e^- -> MnO2 + 2 H2O}
\item \ce{MnO4^- + 2 H2O + 5 e^- -> MnO2 + 4 OH^-}
\item \ce{MnO4^- + 4 H2O + 3 e^- -> MnO2 + 8 OH^-}
\item \ce{MnO4^- + 2 H2O + 3 e^- -> MnO2 + 4 OH^-}
\end{enumerate}
\textit{Bonne réponse : D}\par
\ce{Mn} passe de $+\mathrm{VII}$ à $+\mathrm{IV}$ : $3$ électrons. En milieu basique on n'écrit pas de \ce{H+}. Pour \ce{MnO4^- + 2 H2O + 3 e^- -> MnO2 + 4 OH^-} : oxygène $4+2=6=2+4$ ; hydrogène $4=4$ ; charge $-1+0-3=-4=4\times(-1)$. Tout est conservé.\par
\textit{Pièges :}
\begin{itemize}
\item A: demi-équation écrite en milieu acide (présence de \ce{H+}), interdite en milieu basique.
\item B: produit \ce{MnO2} confondu avec \ce{Mn^{2+}}, d'où $5$ électrons ($+\mathrm{VII}\to+\mathrm{II}$) au lieu de $3$ ; la charge n'est plus équilibrée.
\item C: \ce{H2O} et \ce{OH^-} ajoutés en double : l'oxygène n'est plus conservé ($10$ à droite contre $8$ à gauche).
\end{itemize}
\vspace{8pt}\hrule\vspace{8pt}
\noindent\textbf{CHIM-F11-Q06 --- Prévision d'une réaction spontanée par les potentiels standard}\par
On donne les potentiels standard de réduction (en volt) : \ce{Ag+}/\ce{Ag} : \num{+0.80} ; \ce{Fe^{3+}}/\ce{Fe^{2+}} : \num{+0.77} ; \ce{I2}/\ce{I^-} : \num{+0.54} ; \ce{Cu^{2+}}/\ce{Cu} : \num{+0.34} ; \ce{Zn^{2+}}/\ce{Zn} : \num{-0.76}. Quelle réaction est thermodynamiquement spontanée dans le sens direct (conditions standard) ?
\begin{enumerate}[label=\Alph*)]
\item \ce{Cu^{2+} + 2 Ag -> Cu + 2 Ag+}
\item \ce{2 Fe^{3+} + Cu -> 2 Fe^{2+} + Cu^{2+}}
\item \ce{2 Fe^{2+} + I2 -> 2 Fe^{3+} + 2 I^-}
\item \ce{Zn^{2+} + Cu -> Zn + Cu^{2+}}
\end{enumerate}
\textit{Bonne réponse : B}\par
Une réaction a lieu si l'oxydant appartient au couple de potentiel le plus élevé (règle du $\gamma$, $\Delta E^\circ>0$). Pour \ce{2 Fe^{3+} + Cu -> 2 Fe^{2+} + Cu^{2+}} : \ce{Fe^{3+}} ($+0{,}77$) oxyde \ce{Cu} (couple \ce{Cu^{2+}}/\ce{Cu} à $+0{,}34$), soit $\Delta E^\circ=0{,}77-0{,}34=+0{,}43$~V $>0$.\par
\textit{Pièges :}
\begin{itemize}
\item A: règle du $\gamma$ inversée — \ce{Cu^{2+}} ($+0{,}34$) ne peut pas oxyder \ce{Ag} (couple à $+0{,}80$, situé plus haut).
\item C: on suppose qu'un halogène \ce{I2} oxyde toujours ; or \ce{I2}/\ce{I^-} ($+0{,}54$) est sous \ce{Fe^{3+}}/\ce{Fe^{2+}} ($+0{,}77$), la réaction est impossible.
\item D: \ce{Zn^{2+}} ($-0{,}76$) pris pour un oxydant du cuivre, alors que ce couple est le plus bas de la liste.
\end{itemize}
\vspace{8pt}\hrule\vspace{8pt}
\noindent\textbf{CHIM-F11-Q07 --- Formule de Nernst : potentiel d'une demi-pile}\par
Une électrode d'argent est plongée dans une solution où la concentration en ions \ce{Ag+} est \qty{e-3}{\mole\per\litre}, à \qty{25}{\celsius}. Le potentiel standard $E^\circ_{\mathrm{r}}(\ce{Ag+}/\ce{Ag})$ vaut \qty{+0.80}{\volt}. En utilisant $E=E^\circ_{\mathrm{r}}+\frac{0{,}06}{n}\log[\ce{Ox}]$, quel est le potentiel de cette demi-pile ?
\begin{enumerate}[label=\Alph*)]
\item \qty{0.62}{\volt}
\item \qty{0.71}{\volt}
\item \qty{0.80}{\volt}
\item \qty{0.98}{\volt}
\end{enumerate}
\textit{Bonne réponse : A}\par
Pour \ce{Ag+ + e^- <=> Ag}, $n=1$ et le métal \ce{Ag} a une activité $1$. Nernst : $E=0{,}80+\frac{0{,}06}{1}\log(10^{-3})=0{,}80+0{,}06\times(-3)=0{,}80-0{,}18=0{,}62$~V. Diluer la forme oxydée abaisse le potentiel.\par
\textit{Pièges :}
\begin{itemize}
\item B: $n=2$ utilisé au lieu de $n=1$ pour \ce{Ag+}/\ce{Ag} : $0{,}80+\frac{0{,}06}{2}\times(-3)=0{,}71$~V.
\item C: correction de Nernst oubliée, on garde $E=E^\circ_{\mathrm{r}}=0{,}80$~V.
\item D: erreur de signe sur le logarithme ($\log 10^{-3}$ compté $+3$) : $0{,}80+0{,}18=0{,}98$~V.
\end{itemize}
\vspace{8pt}\hrule\vspace{8pt}
\noindent\textbf{CHIM-F11-Q08 --- Force électromotrice standard d'une pile}\par
Pour la pile Daniell \ce{Zn|Zn^{2+}||Cu^{2+}|Cu} en conditions standard, on donne $E^\circ_{\mathrm{r}}(\ce{Cu^{2+}}/\ce{Cu})=+0{,}34$~V et $E^\circ_{\mathrm{r}}(\ce{Zn^{2+}}/\ce{Zn})=-0{,}76$~V. Que vaut sa force électromotrice standard $\Delta E^\circ$ ?
\begin{enumerate}[label=\Alph*)]
\item \qty{-1.10}{\volt}
\item \qty{0.42}{\volt}
\item \qty{1.10}{\volt}
\item \qty{2.20}{\volt}
\end{enumerate}
\textit{Bonne réponse : C}\par
Le cuivre (potentiel le plus élevé) est la cathode, le zinc l'anode. $\Delta E^\circ=E^\circ_{\mathrm{r}}(\text{cathode})-E^\circ_{\mathrm{r}}(\text{anode})=0{,}34-(-0{,}76)=+1{,}10$~V, valeur positive : la pile débite spontanément.\par
\textit{Pièges :}
\begin{itemize}
\item A: f.é.m. calculée « anode $-$ cathode » : $-0{,}76-0{,}34=-1{,}10$~V.
\item B: erreur de signe, potentiel du zinc compté $+0{,}76$ : $0{,}34-0{,}76=-0{,}42$~V (valeur absolue $0{,}42$).
\item D: $\Delta E^\circ$ multipliée par le nombre d'électrons $n=2$ (confusion avec une grandeur extensive) : $2{,}20$~V.
\end{itemize}
\vspace{8pt}\hrule\vspace{8pt}
\noindent\textbf{CHIM-F11-Q09 --- Fonctionnement d'une pile galvanique (anode, cathode, porteurs)}\par
Une pile associe les couples \ce{Zn^{2+}}/\ce{Zn} et \ce{Ag+}/\ce{Ag}, reliés par un pont salin. Données : $E^\circ_{\mathrm{r}}(\ce{Ag+}/\ce{Ag})$ vaut \qty{+0.80}{\volt} et $E^\circ_{\mathrm{r}}(\ce{Zn^{2+}}/\ce{Zn})$ vaut \qty{-0.76}{\volt}. Parmi les affirmations suivantes, laquelle est \textbf{correcte} ?
\begin{enumerate}[label=\Alph*)]
\item L'électrode d'argent (cathode, pôle $+$) voit sa masse augmenter, car \ce{Ag+ + e^- -> Ag}.
\item Dans le fil, les électrons circulent de l'argent vers le zinc.
\item Le zinc est la cathode, puisqu'il y est le siège d'une oxydation.
\item Dans le pont salin, les cations migrent vers le compartiment du zinc.
\end{enumerate}
\textit{Bonne réponse : A}\par
Le couple \ce{Ag+}/\ce{Ag} ($+0{,}80$) a le potentiel le plus élevé : l'argent est la cathode (pôle $+$, réduction \ce{Ag+ + e^- -> Ag}), son électrode s'épaissit. Le zinc ($-0{,}76$) est l'anode (pôle $-$, oxydation \ce{Zn -> Zn^{2+} + 2 e^-}) et se dissout.\par
\textit{Pièges :}
\begin{itemize}
\item B: sens des électrons inversé — ils vont de l'anode (\ce{Zn}) vers la cathode (\ce{Ag}) ; c'est le courant conventionnel qui va en sens inverse.
\item C: confusion anode/cathode — l'oxydation a lieu à l'\textbf{anode} (pôle $-$), le zinc n'est donc pas la cathode.
\item D: migration ionique inversée — les cations migrent vers la cathode (\ce{Ag}), les anions vers l'anode (\ce{Zn}).
\end{itemize}
\vspace{8pt}\hrule\vspace{8pt}
\noindent\textbf{CHIM-F11-Q10 --- Loi de Faraday : masse déposée par électrolyse}\par
On électrolyse une solution de sulfate de cuivre \ce{CuSO4} avec une intensité $I=\qty{2.0}{A}$ pendant $t=\qty{45}{min}$. À la cathode : \ce{Cu^{2+} + 2 e^- -> Cu}. Données : $M(\ce{Cu})=\qty{63.5}{\gram\per\mole}$, $F=\qty{96500}{\coulomb\per\mole}$. Quelle masse de cuivre se dépose ?
\begin{enumerate}[label=\Alph*)]
\item \qty{0.030}{\gram}
\item \qty{1.78}{\gram}
\item \qty{3.55}{\gram}
\item \qty{7.11}{\gram}
\end{enumerate}
\textit{Bonne réponse : B}\par
Loi de Faraday : $m=\dfrac{M\,I\,t}{n\,F}$ avec $n=2$ et $t=45\times60=\qty{2700}{s}$. $m=\dfrac{63{,}5\times2{,}0\times2700}{2\times96500}=\dfrac{342900}{193000}\simeq 1{,}78$~g.\par
\textit{Pièges :}
\begin{itemize}
\item A: durée laissée en minutes ($t=45$ au lieu de \qty{2700}{s}) : $m\simeq 0{,}030$~g.
\item C: on oublie que \ce{Cu^{2+}} capte $2$ électrons (on prend $n=1$) : $m=\dfrac{342900}{96500}\simeq 3{,}55$~g.
\item D: le facteur $n$ placé au numérateur ($m=\frac{M\,I\,t\,n}{F}$) au lieu du dénominateur : $m\simeq 7{,}11$~g.
\end{itemize}
\vspace{8pt}\hrule\vspace{8pt}
\end{document}
